\documentclass{article}
\usepackage{amsmath, amsthm}

\newtheorem{proposicion}{Proposición}
\newtheorem{corolario}{Corolario}

\begin{document}

\section*{Criterios de Convergencia de Series}

%--------------------------------------------------
\begin{proposicion}[Criterio de acotación]
Una serie de términos no negativos $\sum_{n=1}^{\infty} a_n$ converge si y sólo si
sus sumas parciales están acotadas.
\end{proposicion}

%--------------------------------------------------
\begin{proposicion}[Criterio de Cauchy para series]
Una serie $\sum_{n=1}^{\infty} a_n$ converge si y sólo si para todo $\varepsilon > 0$
existe $N \in \mathbb{N}$ tal que
\[
  \left|\sum_{i=n}^{m} a_i\right| < \varepsilon
  \quad \text{para todo } m \geq n > N.
\]
\end{proposicion}

%--------------------------------------------------
\begin{proposicion}[Condición del $n$-ésimo término]
Si $\sum_{n=1}^{\infty} a_n$ converge, entonces $a_n \to 0$.
\end{proposicion}

%--------------------------------------------------
\begin{proposicion}[Criterio de comparación directa]
Supongamos que existe $N \in \mathbb{N}$ tal que
\[
  0 \leq a_n \leq b_n \quad \text{para todo } n > N.
\]
\begin{enumerate}
  \item[(a)] Si $\sum_{n=1}^{\infty} b_n$ converge, entonces $\sum_{n=1}^{\infty} a_n$ converge.
  \item[(b)] Si $\sum_{n=1}^{\infty} a_n$ diverge, entonces $\sum_{n=1}^{\infty} b_n$ diverge.
\end{enumerate}
\end{proposicion}

%--------------------------------------------------
\begin{proposicion}[Criterio de comparación al límite del cociente]
Si $a_n, b_n > 0$ y $\lim_{n\to\infty} a_n/b_n$ es un número real distinto de cero,
entonces
\[
  \sum_{n=1}^{\infty} a_n \text{ converge} \iff \sum_{n=1}^{\infty} b_n \text{ converge}.
\]
\end{proposicion}

%--------------------------------------------------
\begin{proposicion}[Criterio de la integral]
Supongamos que $f:[1,\infty)\to(0,\infty)$ es monótona decreciente. Entonces
\[
  \sum_{n=1}^{\infty} f(n) \text{ converge} \iff \int_1^{\infty} f \text{ converge}.
\]
\end{proposicion}

%--------------------------------------------------
\begin{proposicion}[$p$-series]
Sea $p > 0$.
\begin{enumerate}
  \item[(a)] Si $p > 1$, entonces $\displaystyle\sum_{n=1}^{\infty} \frac{1}{n^p}$ converge.
  \item[(b)] Si $0 < p \leq 1$, entonces $\displaystyle\sum_{n=1}^{\infty} \frac{1}{n^p}$ diverge.
\end{enumerate}
\end{proposicion}

%--------------------------------------------------
\begin{proposicion}[Criterio de la raíz]
Consideremos una serie $\sum_{n=1}^{\infty} a_n$ de términos no negativos y definamos
\[
  \alpha = \lim_{n\to\infty} \sqrt[n]{a_n}.
\]
\begin{enumerate}
  \item[(a)] Si $\alpha < 1$, entonces $\sum_{n=1}^{\infty} a_n$ converge.
  \item[(b)] Si $\alpha > 1$, entonces $\sum_{n=1}^{\infty} a_n$ diverge.
  \item[(c)] Si $\alpha = 1$, entonces $\sum_{n=1}^{\infty} a_n$ puede converger o diverger.
\end{enumerate}
\end{proposicion}

%--------------------------------------------------
\begin{proposicion}[Criterio del cociente]
Consideremos una serie $\sum_{n=1}^{\infty} a_n$ de términos positivos y definamos
\[
  \alpha = \lim_{n\to\infty} \frac{a_{n+1}}{a_n}.
\]
\begin{enumerate}
  \item[(a)] Si $\alpha < 1$, entonces $\sum_{n=1}^{\infty} a_n$ converge.
  \item[(b)] Si $\alpha > 1$, entonces $\sum_{n=1}^{\infty} a_n$ diverge.
  \item[(c)] Si $\alpha = 1$, entonces $\sum_{n=1}^{\infty} a_n$ puede converger o diverger.
\end{enumerate}
\end{proposicion}

%--------------------------------------------------
\begin{corolario}[Criterio de Raabe]
Consideremos una serie $\sum_{n=1}^{\infty} a_n$ de términos positivos y definamos
\[
  \alpha = \lim_{n\to\infty} n\!\left(1 - \frac{a_{n+1}}{a_n}\right).
\]
\begin{enumerate}
  \item[(a)] Si $\alpha > 1$, entonces $\sum_{n=1}^{\infty} a_n$ converge.
  \item[(b)] Si $\alpha < 1$, entonces $\sum_{n=1}^{\infty} a_n$ diverge.
  \item[(c)] Si $\alpha = 1$, entonces $\sum_{n=1}^{\infty} a_n$ puede converger o diverger.
\end{enumerate}
\end{corolario}

%--------------------------------------------------
\begin{proposicion}[Criterio de Leibniz]
Si $(a_n)$ es una sucesión monótona decreciente que converge a cero, entonces
las series alternantes
\[
  \sum_{n=1}^{\infty} (-1)^n a_n \qquad \text{y} \qquad \sum_{n=1}^{\infty} (-1)^{n+1} a_n
\]
convergen.
\end{proposicion}

%--------------------------------------------------
\begin{proposicion}[Convergencia absoluta implica convergencia]
Si $\sum_{n=1}^{\infty} a_n$ converge absolutamente, entonces converge y
\[
  \left|\sum_{n=1}^{\infty} a_n\right| \leq \sum_{n=1}^{\infty} |a_n|.
\]
\end{proposicion}

%--------------------------------------------------
\section*{Criterio para Series de Funciones}

\begin{proposicion}[Criterio $M$ de Weierstrass]
Sean $(f_n)$ una sucesión de funciones definidas en $D$ y $(M_n)$ una sucesión
de números reales tales que
\[
  |f_n(x)| \leq M_n \quad \text{para todo } x \in D.
\]
Si $\sum_{n=1}^{\infty} M_n$ converge, entonces $\sum_{n=1}^{\infty} f_n$ converge uniformemente en $D$.
\end{proposicion}

\end{document}